% ============================================================================
% Fichier  : partie1/P01_C02_principes.tex
% Titre    : Les Principes du Traitement
% Auteur   : MINKA MI NGUIDJOÏ Thierry Emmanuel
% Version  : 1.0 -- novembre 2026
% Corrigés : corriges/Corr_C02.tex
% ============================================================================
% Sources : loi n° 2024/017, art. 5 (consentement), 6 à 18, 21, 36, 50
% (texte officiel) ; ISO/IEC 29100 (principes, paraphrasés) ; RGPD art. 6
% (comparaison).
% Faits nouveaux (bible) : formulaire d'ouverture de compte de 2016 ; comptes
% clos depuis 2012 ; prospects collectés par les agents ; 14 juillet 2026.
% ============================================================================

\chapter{Les Principes du Traitement}
\label{chap:principes}
\chaptermark{Les principes du traitement}

\epigraph{%
    « Une case cochée n'est pas un consentement.\\
    C'est, au mieux, un indice. »%
}{Principe du cours}

\niveauChapitre{Fondamental}{%
    Chapitre~\ref{chap:donnees}
}

\begin{boxobjectifs}
\begin{itemize}
    \item Énoncer les principes des articles~6 à~18 de la loi n°~2024/017
          et les traduire en questions d'audit.
    \item Apprécier la validité d'un consentement au regard des
          articles~5 et~9.
    \item Identifier les trois cas où l'on peut se passer du consentement,
          et les comparer aux bases légales du droit européen.
    \item Vérifier la finalité, l'exactitude et la durée de conservation
          d'un traitement.
    \item Contrôler l'information fournie aux personnes selon
          l'article~21.
\end{itemize}
\end{boxobjectifs}

\bigskip

% ============================================================================
\section*{Avant de commencer : « ils ont tous signé »}
% ============================================================================

Bafoussam, mardi 14~juillet~2026. Le directeur commercial de NYANGA CRÉDIT
pose sur la table de Mme~BELLO une liasse de formulaires d'ouverture de
compte. « Ils ont tous signé. » Le formulaire, inchangé depuis 2016, tient
sur une page. En bas, au-dessus de la signature, une case pré-imprimée~:
« J'accepte les conditions générales de NYANGA CRÉDIT. » Les conditions
générales, huit pages, ne sont pas remises au client~; elles sont
« disponibles en agence ».

L'auditrice note trois questions. Les clients ont-ils consenti à quelque
chose de précis~? Les cautions, qui ne signent pas ce formulaire, ont-elles
consenti à quoi que ce soit~? Et les prospects, dont les agents notent le
nom et le téléphone sur un cahier lors des tournées de marché, savent-ils
seulement qu'ils figurent dans un fichier~?

\begin{boxalert}
\textbf{Violation typique.} Tenir une signature au bas de conditions
générales pour un consentement au traitement des données.

\smallskip
\textbf{Règle~:} le consentement exigé par la loi est préalable, libre,
éclairé, spécifique et univoque. Il porte sur un traitement identifié, après
une information claire~; il se prouve, et il peut être retiré.
\end{boxalert}

% ============================================================================
\section{Six Principes, Six Questions d'Audit}
\label{sec:c02-principes}
% ============================================================================

Le chapitre~I du titre~II de la loi pose les principes du traitement. Chacun
se traduit en une question fermée que l'auditeur peut coter.

\begin{tableaucours}{L{3.0cm}cL{8.0cm}}{%
    \tableauHeader{Principe} &
    \tableauHeader{Article} &
    \tableauHeader{Question d'audit}
}
Respect de la vie privée & 6 & Le traitement porte-t-il atteinte à la vie
privée au-delà de ce que sa finalité exige~? \\
Confidentialité & 7, 15 & Les données sont-elles protégées contre la
divulgation, y compris en transmission~? \\
Consentement & 9 & Un consentement valable, ou l'une des trois dérogations,
fonde-t-il le traitement~? \\
Licéité et loyauté & 10 & Le traitement est-il permis, et les personnes
ont-elles été traitées loyalement~? \\
Finalité & 11 & La finalité est-elle déterminée, explicite et légitime, et
l'usage ultérieur compatible~? \\
Exactitude & 12 & Les données sont-elles complètes, fiables, tenues à jour,
et les erreurs corrigées~? \\
Conservation limitée & 13 & La durée de conservation est-elle fixée et
respectée~? \\
Information & 14, 21 & Les personnes reçoivent-elles, à la collecte, les
informations de l'article~21~? \\
Non-discrimination & 17 & Le traitement est-il indifférent à la classe
sociale, à l'origine, aux opinions~? \\
Mineurs & 18 & Le traitement des données de mineurs est-il adéquat,
pertinent et limité~? \\
\end{tableaucours}
\vspace{-0.8em}
\begin{center}{\small\itshape Tableau~2.1. Les principes de la loi et leurs questions d'audit}\end{center}

\begin{boxinfo}
\textbf{La minimisation n'est pas un principe général de la loi.} Le mot
n'apparaît qu'à l'article~36, pour l'interconnexion des fichiers, et l'idée
de données « adéquates, pertinentes et limitées » n'est formulée qu'à
l'article~18, pour les mineurs. L'auditeur ne peut donc pas coter un écart à
un principe général de minimisation au titre de la loi. Il peut en revanche
s'appuyer sur la finalité (article~11)~: une donnée collectée sans lien avec
la finalité annoncée n'a pas de fondement. Les normes, elles, posent la
minimisation comme principe~\cite{iso29100}~; un audit contre l'ISO/IEC~27701
peut donc la retenir comme critère, à condition de le dire.
\end{boxinfo}

% ============================================================================
\section{Le Consentement}
\label{sec:c02-consentement}
% ============================================================================

\begin{boxdef}
\termecle{Consentement}~: selon l'article~5, « toute manifestation de volonté
expresse, non équivoque, non contrainte, libre et éclairée par une
information claire, précise et complète, par laquelle la personne
concernée, ou son représentant légal, accepte que ses données à caractère
personnel fassent l'objet d'un traitement ».
\end{boxdef}

\begin{boxjuris}
\textbf{\loicam{n°~2024/017}, article~9}~\cite{loi2024017}~:
\textit{« (1) Le traitement des données à caractère personnel est subordonné
au consentement préalable, libre, éclairé, spécifique et univoque de la
personne concernée. (2) Toutefois, il peut être dérogé à cette exigence du
consentement préalable lorsque le traitement est nécessaire~: au respect
d'une obligation légale à laquelle le responsable du traitement est soumis~;
à l'exécution d'une mission d'intérêt public ou relevant de l'exercice de
l'Autorité de protection des données à caractère personnel, dont est
investi le responsable du traitement ou le tiers auquel les données sont
communiquées~; à la préservation de la santé de la personne concernée ou
d'une autre personne physique. (3) Le consentement d'une personne mineure de
dix-huit (18) ans n'est valable que s'il est appuyé par celui de ses père
et mère, ou de son représentant légal. »}

\smallskip
\textbf{Portée.} Le consentement est la règle~; les dérogations sont trois,
et elles sont étroites. L'article~50 le confirme de façon négative~:
« le traitement des données à caractère personnel sans le consentement
préalable de la personne concernée est interdit ».
\end{boxjuris}

\subsection{Les cinq qualités d'un consentement}

Pour coter \OUI{} la proposition « le traitement repose sur un consentement
valable », l'auditeur vérifie cinq qualités, chacune par une pièce.

\begin{enumerate}
    \item \textbf{Préalable}~: recueilli avant le traitement. Pièce~: la date
          du recueil, comparée à la date de la première donnée.
    \item \textbf{Libre}~: sans contrainte. Un consentement exigé comme
          condition d'un service qui n'en a pas besoin n'est pas libre.
    \item \textbf{Éclairé}~: après une information claire, précise et
          complète. Pièce~: le texte remis au moment du recueil.
    \item \textbf{Spécifique}~: pour une finalité identifiée. Un
          consentement global à « toutes les opérations » ne l'est pas.
    \item \textbf{Univoque}~: un acte positif, sans ambiguïté. Une case
          pré-cochée ou un silence ne le sont pas.
\end{enumerate}

À ces cinq qualités s'ajoute la \textbf{preuve}~: c'est au responsable de
montrer qu'il a recueilli le consentement. Un consentement dont on ne
retrouve pas la trace est, pour l'auditeur, un consentement \NEK{}, et
souvent un \NON{} si la procédure de recueil n'existait pas.

\subsection{Les trois dérogations, comparées au droit européen}

Le droit européen connaît six bases légales~: le consentement, le contrat,
l'obligation légale, les intérêts vitaux, la mission d'intérêt public et
l'intérêt légitime~\cite{ue2016rgpd}. La loi camerounaise n'en retient que
quatre~: le consentement et trois dérogations. Elle ne connaît ni
l'\textbf{exécution d'un contrat}, ni l'\textbf{intérêt légitime}. Le
Chapitre~\ref{chap:rgpd} reprend les six bases une à une et montre ce que
cet écart change pour une organisation soumise aux deux textes.

\begin{tableaucours}{L{4.4cm}L{3.6cm}L{4.0cm}}{%
    \tableauHeader{Situation à NYANGA} &
    \tableauHeader{Fondement possible} &
    \tableauHeader{Commentaire}
}
Vérifier l'identité d'un client à l'ouverture & Obligation légale & Si un
texte de lutte contre le blanchiment l'impose \\
Tenir le compte et exécuter les opérations & Consentement & Pas de base
« contrat » dans la loi \\
Vidéosurveillance des agences & Consentement~? & Pas d'intérêt légitime~:
fondement fragile \\
Relancer un débiteur par un prestataire & Consentement & À recueillir
expressément \\
Démarcher un prospect au marché & Consentement & Préalable à l'inscription
au cahier \\
Transmettre les fichiers à l'expert judiciaire & Obligation légale &
Réquisition judiciaire \\
\end{tableaucours}
\vspace{-0.8em}
\begin{center}{\small\itshape Tableau~2.2. Fondements des traitements de NYANGA CRÉDIT}\end{center}

L'absence de l'intérêt légitime a des conséquences lourdes. La
vidéosurveillance, la sécurité informatique ou la prévention de la fraude,
qui reposent en Europe sur l'intérêt légitime, doivent trouver au Cameroun
un consentement, une obligation légale ou une mission d'intérêt public.
L'auditeur ne comble pas ce vide en important le droit européen~: il
signale le traitement dont le fondement est incertain, le cote \NEK{}, et
recommande d'obtenir l'autorisation de l'Autorité, qui en fixera les
conditions (article~19).

% ============================================================================
\section{Finalité, Exactitude, Conservation}
\label{sec:c02-finalite}
% ============================================================================

\begin{boxjuris}
\textbf{\loicam{n°~2024/017}, article~11}~\cite{loi2024017}~:
\textit{« Les données à caractère personnel doivent être collectées pour des
finalités déterminées, explicites et légitimes. Elles ne peuvent pas être
ultérieurement traitées de manière incompatible avec ces finalités. »}
\end{boxjuris}

Le détournement de finalité est l'écart le plus fréquent, et l'un des plus
graves~: l'article~67 en fait un délit. Le fichier des débiteurs, constitué
pour le recouvrement, devient incompatible s'il sert à démarcher les mêmes
personnes pour un nouveau crédit~; les images des caméras, prises pour la
sécurité, le deviennent si elles servent à contrôler les horaires des
agents.

L'article~12 impose des données « complètes, fiables et tenues à jour » et
oblige à corriger ou effacer les données inexactes. L'auditeur le teste par
sondage~: il tire des dossiers au hasard et compare les données aux pièces
justificatives.

\begin{boxjuris}
\textbf{\loicam{n°~2024/017}, article~13}~\cite{loi2024017}~:
\textit{« (1) Les données à caractère personnel doivent être conservées
pendant une durée qui n'excède pas la période nécessaire à leur traitement.
(2) La période de conservation des données à caractère personnel visée à
l'alinéa 1 ci-dessus est précisée dans l'autorisation de traitement des
données et suivant les modalités fixées par voie réglementaire. (3) Au-delà
de la période visée à l'alinéa 1 ci-dessus, les données à caractère
personnel ne peuvent être conservées qu'en vue de répondre spécifiquement à
un traitement à des fins historiques, statistiques ou de recherche
conformément aux dispositions légales en vigueur. »}

\smallskip
\textbf{Portée.} La durée se fixe dans l'autorisation~; tant que celle-ci
n'existe pas, l'organisation doit fixer elle-même une durée justifiée, et
l'auditeur vérifie qu'elle l'applique. L'article~28 renvoie aussi au
référentiel de l'Autorité.
\end{boxjuris}

% ============================================================================
\section{L'Information des Personnes}
\label{sec:c02-information}
% ============================================================================

L'article~14 exige une information « claire et transparente », et
l'article~21 en fixe le contenu~: elle doit être fournie « au plus tard, lors
de la collecte et quels que soient les moyens employés ».

\begin{tableaucours}{L{6.0cm}L{6.2cm}}{%
    \tableauHeader{Élément exigé (article~21)} &
    \tableauHeader{Ce que l'auditeur cherche}
}
Identité du responsable & Nom et coordonnées sur le formulaire \\
Finalités déterminées & Une phrase par finalité, sans formule globale \\
Catégories de données & La liste des données collectées \\
Destinataires & Les catégories de destinataires, prestataires compris \\
Possibilité de refuser de figurer au fichier & Une mention explicite \\
Droits (accès, rectification, effacement, opposition, limitation,
portabilité, décision automatisée, directives post-mortem, réclamation) &
La liste des droits et la manière de les exercer \\
Durée de conservation & Une durée ou un critère \\
Transfert éventuel vers un État étranger & La mention du pays ou de la
région \\
\end{tableaucours}
\vspace{-0.8em}
\begin{center}{\small\itshape Tableau~2.3. Le contenu de l'information (article~21)}\end{center}

La liste est longue, et elle se contrôle mécaniquement~: l'auditeur pose le
formulaire à côté du tableau et coche. L'information peut être donnée en
deux niveaux, une mention courte sur le formulaire et une notice complète
remise ou accessible, pourvu que la personne la reçoive vraiment au moment
de la collecte. « Disponible en agence » ne suffit pas.

% ============================================================================
\section{Les Principes à NYANGA}
\label{sec:c02-nyanga}
% ============================================================================

\begin{boxscenario}{Mission NYANGA. Le formulaire de 2016 et les autres, 14 juillet 2026}
\textbf{Constatations.}
\begin{itemize}
    \item Le formulaire d'ouverture de compte porte une seule case,
          pré-imprimée, d'acceptation des conditions générales~; aucune
          finalité n'y figure. Les conditions générales ne sont pas remises.
    \item Les cautions signent un acte de cautionnement qui ne mentionne
          pas le traitement de leurs données.
    \item Les cahiers de prospection de trois agences tirées au hasard
          contiennent 1\,214 noms et numéros de téléphone, sans aucune
          mention d'information ni de consentement.
    \item La base des comptes contient 61\,380 comptes clos, dont 22\,950
          clos avant 2016, conservés avec toutes leurs données.
    \item Un sondage de 40 dossiers actifs trouve 9 numéros de téléphone
          erronés et 6 adresses périmées.
\end{itemize}

\textbf{Analyse} (confiance élevée). Le consentement des clients n'est ni
spécifique, ni éclairé~: il ne fonde pas les traitements au sens de
l'article~9. Une partie des données de compte peut relever d'une obligation
légale d'identification, que l'audit n'a pas encore établie. Les cautions et les prospects n'ont
consenti à rien. Les comptes clos depuis dix ans et plus n'ont plus de
finalité identifiée.
\end{boxscenario}

\begin{boxlecture}{Les principes à NYANGA CRÉDIT}
\begin{tableaucours}{L{7.6cm}cL{3.4cm}}{%
    \tableauHeader{Proposition} &
    \tableauHeader{Cote} &
    \tableauHeader{Pièce}
}
Les clients ont donné un consentement valable & \NON{} & Formulaire 2016 \\
Les cautions ont consenti au traitement & \NON{} & Acte de cautionnement \\
Les prospects ont consenti au traitement & \NON{} & Cahiers de prospection \\
Une obligation légale fonde l'identification des clients & \NEK{} & Texte non produit \\
Une durée de conservation est fixée & \NON{} & Aucune règle écrite \\
Les données des comptes actifs sont exactes & \NON{} & Sondage, 15 erreurs
sur 40 \\
L'information de l'article~21 est fournie & \NON{} & Formulaire 2016 \\
La vidéosurveillance repose sur un fondement valable & \NEK{} & Pas de
consentement, pas d'autorisation \\
\end{tableaucours}

\medskip
\textbf{Décompte.} Zéro \OUI{}, six \NON{}, deux \NEK{} sur huit.
\hfill\lecture{0}{0{,}25}{0{,}81}

\medskip
\textbf{Lecture.} Le domaine des principes est clair~: $U$ est faible et
les écarts sont établis. L'entropie est basse parce que les réponses se
concentrent sur \NON{}. Une conclusion de non-conformité peut être donnée
dès maintenant, avec deux points à instruire.
\end{boxlecture}

% ============================================================================
\section{Ce que les Principes demandent à chaque Métier}
\label{sec:c02-metiers}
% ============================================================================

\begin{boxmetier}{Auditeur de la protection des données}
Pour chaque traitement, demander le fondement, puis la pièce qui le prouve.
Tester les cinq qualités du consentement une par une, et ne jamais
importer une base légale européenne que la loi ne connaît pas.
\end{boxmetier}

\begin{boxmetier}{Juriste conformité}
Réécrire les formulaires~: une finalité par phrase, une case non cochée par
consentement distinct, la notice de l'article~21 remise au moment de la
collecte. Identifier les textes qui imposent une obligation légale.
\end{boxmetier}

\begin{boxmetier}{Délégué à la protection des données}
Tenir, pour chaque traitement fondé sur le consentement, la preuve de son
recueil et la procédure de retrait.
\end{boxmetier}

\begin{boxresume}
\begin{itemize}
    \item Les principes des articles~6 à~18 se traduisent chacun en une
          question d'audit fermée.
    \item Le consentement est la règle~; il doit être préalable, libre,
          éclairé, spécifique et univoque, et sa preuve incombe au
          responsable.
    \item La loi ne connaît que trois dérogations au consentement~; ni le
          contrat ni l'intérêt légitime n'en font partie.
    \item La durée de conservation se fixe dans l'autorisation~; à défaut,
          l'organisation doit la fixer et l'appliquer.
    \item À NYANGA CRÉDIT, six écarts sont établis sur huit propositions~:
          consentements, durée, exactitude, information.
\end{itemize}
\end{boxresume}

\begin{boxinfo}
\textbf{Points de contrôle.} Vous avez compris ce chapitre si vous savez~:
\begin{itemize}
    \item tester les cinq qualités d'un consentement sur un formulaire~;
    \item expliquer pourquoi la vidéosurveillance pose un problème de
          fondement au Cameroun~;
    \item contrôler une notice d'information contre l'article~21.
\end{itemize}
\end{boxinfo}

\begin{boxexo}[title={\ding{45}\quad Exercice~2.1. Valable ou non \niveaubase}]
Le consentement est-il valable~? (a) une case pré-cochée « J'accepte de
recevoir des offres »~; (b) un client qui dicte son accord par téléphone,
enregistré~; (c) un consentement exigé pour ouvrir un compte d'épargne à
une utilisation de ses données pour la prospection~; (d) un élève de
16~ans qui s'inscrit seul à un service de microcrédit scolaire.
\end{boxexo}

\begin{boxexo}[title={\ding{45}\quad Exercice~2.2. Réécrire le formulaire \niveaumoyen}]
Réécrivez le formulaire d'ouverture de compte de NYANGA CRÉDIT pour qu'il
respecte les articles~9 et~21. Distinguez ce qui relève d'une obligation
légale et ce qui exige un consentement.
\end{boxexo}

\begin{boxexo}[title={\ding{45}\quad Exercice~2.3. Les comptes clos \niveaumoyen}]
\begin{enumerate}[(a)]
    \item Proposez une règle de conservation pour les comptes clos, en
          justifiant chaque durée.
    \item Que faire des 22\,950 comptes clos avant 2016~?
    \item Que faire des données de M.~TAGNE si son compte était clos~?
\end{enumerate}
\end{boxexo}

\begin{boxexo}[title={\ding{45}\quad Exercice~2.4. Sans intérêt légitime \niveauexpert}]
NOUN PAY veut analyser les transactions de ses clients pour détecter la
fraude. En Europe, elle invoquerait l'intérêt légitime. Quels fondements
peut-elle invoquer au Cameroun~? Quelles pièces l'auditeur demandera-t-il,
et comment cotez-vous la proposition « la détection de fraude repose sur un
fondement valable » si aucune autorisation n'a été délivrée~?
\end{boxexo}

\bigskip

\begin{boxinfo}
\textbf{Pour aller plus loin.}
\begin{itemize}
    \item L'ISO/IEC~29100~\cite{iso29100}, pour les principes de
          consentement, de limitation de la collecte et de minimisation.
    \item Le RGPD~\cite{ue2016rgpd}, article~6, pour comparer les bases
          légales.
\end{itemize}
\end{boxinfo}

\begin{boxinfo}
\textbf{Transition.} Les écarts de NYANGA commencent à se dessiner, mais
une liste de défauts n'est pas un audit. Le Chapitre~\ref{chap:audit}
présente la démarche qui transforme des constats en une conclusion que l'on peut défendre.
\end{boxinfo}
